% 演習問題集 — 第4章
% 章はこのファイル名によって決まる。

\begin{exo}[経路に沿った微分形式の積分]{integrale-forme-differentielle}
\keywords{微分形式, 完全微分, シュワルツの判定条件, 線積分}

二つの微分形式
\[
\omega_1 = y\,\mathrm{d}x + x\,\mathrm{d}y,
\qquad
\omega_2 = y\,\mathrm{d}x,
\]
と、原点 $O(0,0)$ から点 $M(1,1)$ へ至る三つの経路を考える。
$\gamma_1$ は水平線分 $O\to(1,0)$ と、それに続く鉛直線分 $(1,0)\to M$、
$\gamma_2$ は鉛直線分 $O\to(0,1)$ と、それに続く水平線分 $(0,1)\to M$、
$\gamma_3$ は $O$ と $M$ を直接結ぶ直線分である。

\begin{enumerate}
\item $\omega_1 = \mathrm{d}f$ を満たす関数 $f$ を示し、$\omega_1$ が完全形式であることを証明せよ。
\item 各線分をパラメータ表示して $\int_{\gamma_1}\omega_1$ と $\int_{\gamma_2}\omega_1$ を計算せよ。また、計算をせずに同じ結果を導け。
\item $\omega_2$ にシュワルツの判定条件を適用せよ。何が結論できるか。
\item $\int_{\gamma_1}\omega_2$、$\int_{\gamma_2}\omega_2$、$\int_{\gamma_3}\omega_2$ を計算し、考察せよ。
\end{enumerate}

\begin{indication}
水平線分では $\mathrm{d}y = 0$、鉛直線分では $\mathrm{d}x = 0$ であり、
各積分は通常の一変数積分に帰着する。
\end{indication}

\begin{solution}
\textbf{問1。} $f(x,y) = xy$ とすればよい。実際、$\partial f/\partial x = y$、$\partial f/\partial y = x$ である。

\textbf{問2。} $t\mapsto\bigl(x(t),y(t)\bigr)$ とパラメータ表示された曲線に対して
\[
\int_\gamma\omega_1
= \int \left[y(t)x'(t)+x(t)y'(t)\right],\mathrm{d}t.
\]

経路 $\gamma_1$ は二つの線分の和である。第一線分では $x(t)=t$、$y(t)=0$、
$t\in[0,1]$ なので、$\mathrm{d}x=\mathrm{d}t$、$\mathrm{d}y=0$ である。
第二線分では $x(t)=1$、$y(t)=t$、同じく $t\in[0,1]$ なので、
$\mathrm{d}x=0$、$\mathrm{d}y=\mathrm{d}t$ である。したがって
\[
\begin{aligned}
I_1
= \int_{\gamma_1}\omega_1
&= \int_0^1\left(0\times 1+t\times 0\right)\,\mathrm{d}t
 + \int_0^1\left(t\times 0+1\times 1\right)\,\mathrm{d}t \\
&= 0+\int_0^1\mathrm{d}t = 1.
\end{aligned}
\]

$\gamma_2$ では、第一線分を $x(t)=0$、$y(t)=t$、第二線分を
$x(t)=t$、$y(t)=1$ とパラメータ表示する。いずれも $t\in[0,1]$ である。よって
\[
\begin{aligned}
I_2
= \int_{\gamma_2}\omega_1
&= \int_0^1\left(t\times 0+0\times 1\right)\,\mathrm{d}t
 + \int_0^1\left(1\times 1+t\times 0\right)\,\mathrm{d}t \\
&= 0+\int_0^1\mathrm{d}t = 1.
\end{aligned}
\]
二つの値は一致する。これは計算せずとも分かる。$\omega_1=\mathrm{d}f$、
$f(x,y)=xy$ なので、その積分は端点だけに依存し、
\[
\int_\gamma\omega_1=f(M)-f(O)=f(1,1)-f(0,0)=1.
\]

\textbf{問3。} $\omega_2=A(x,y)\,\mathrm{d}x+B(x,y)\,\mathrm{d}y$ と書く。
ここで $A(x,y)=y$、$B(x,y)=0$ である。形式が完全なら、シュワルツの判定条件から
\[
\frac{\partial A}{\partial y}=\frac{\partial B}{\partial x}
\]
でなければならない。しかし $\partial A/\partial y=1$ に対し、
$\partial B/\partial x=0$ である。交差偏微分が異なるので $\omega_2$ は完全形式ではない。
したがって、その積分は $O$ から $M$ までの経路に依存し得る。次の計算でこれを確かめる。

\textbf{問4。} $\omega_2=y\,\mathrm{d}x$ なので、$y$ がゼロでない値にあるときの
$x$ の変化だけが積分に寄与する。先のパラメータ表示を用いると、$\gamma_1$ では
\[
\begin{aligned}
J_1
=\int_{\gamma_1}\omega_2
&=\int_0^1 0\times 1\,\mathrm{d}t
  +\int_0^1 t\times 0\,\mathrm{d}t \\
&=0.
\end{aligned}
\]
$\gamma_2$ の第一線分では $x(t)=0$、したがって $\mathrm{d}x=0$ である。
第二線分では $x(t)=t$、$y(t)=1$、$\mathrm{d}x=\mathrm{d}t$ なので
\[
\begin{aligned}
J_2
=\int_{\gamma_2}\omega_2
&=\int_0^1 t\times 0\,\mathrm{d}t
  +\int_0^1 1\times 1\,\mathrm{d}t \\
&=1.
\end{aligned}
\]
最後に、対角線分 $\gamma_3$ は
\[
x(t)=t,\qquad y(t)=t,\qquad t\in[0,1]
\]
と明示的にパラメータ表示できる。したがって $\mathrm{d}x=\mathrm{d}t$、
$\mathrm{d}y=\mathrm{d}t$ であり、
\[
J_3
=\int_{\gamma_3}\omega_2
=\int_0^1 y(t)x'(t)\,\mathrm{d}t
=\int_0^1 t\,\mathrm{d}t
=\frac12.
\]
三経路の端点は同じだが、$J_1=0$、$J_2=1$、$J_3=1/2$ と異なる値を与える。
これこそ完全でない微分形式に特徴的な経路依存性である。
\end{solution}
\end{exo}

\begin{exo}[同じ内部エネルギー変化、三つの移動形態]{meme-delta-u-trois-transferts}
\keywords{熱力学第一法則, 内部エネルギー, 熱, 力学的仕事, 電気的仕事, 熱量測定, 状態関数}

液体、その熱量計、浸された小さな抵抗器が、巨視的に静止した閉じた系を構成する。
容器は剛体で、系の全熱容量は一定とする：
\[
C=2{,}00\ \mathrm{kJ\,K^{-1}}.
\]
系を、$T_A=20{,}0\,^\circ\mathrm{C}$ の平衡状態 $A$ から
$T_B=25{,}0\,^\circ\mathrm{C}$ の平衡状態 $B$ に移したい。次を仮定する：
\[
U(B)-U(A)=C(T_B-T_A).
\]
巨視的な運動エネルギーと位置エネルギーの変化はゼロである。
外部への損失はすべて無視する。

同じ初期状態 $A$ から、次の三つの手順をそれぞれ独立に行う：
\begin{itemize}
\item[手順1] 系を高温熱源と熱接触させるが、系には仕事を一切与えない。
\item[手順2] 壁を断熱する。質量 $m$ のおもりを高さ $h=5{,}00$ m だけ落下させ、その力で羽根車が液体をかき混ぜる。最後には羽根車と液体は静止し、おもりが失った力学的エネルギーはすべて系に伝わる。$g=9{,}81\ \mathrm{m\,s^{-2}}$ とする。
\item[手順3] 系は熱 $Q_3=4{,}00$ kJ を受け取り、不足分のエネルギーは電気的に抵抗器へ供給される。受け取る電力は一定で $\mathcal P=300$ W である。
\end{itemize}

\begin{enumerate}
\item 三手順に共通する内部エネルギー変化 $\Delta U=U(B)-U(A)$ を求めよ。
\item 最初の二手順のそれぞれについて、系が受け取る熱 $Q$ と仕事 $W$ を求めよ。第二手順では必要な質量 $m$ を求めよ。
\item 第三手順について、受け取る電気的仕事と抵抗器への通電時間を求めよ。
\item 三つの過程は同じ初期状態と最終状態をもつ。それにもかかわらず $Q$ と $W$ が異なり得る理由を説明せよ。最終状態 $B$ の系は「熱」や「仕事」を含むのか。
\end{enumerate}

\begin{indication}
銀行家の符号規約で $\Delta U=Q+W$ を適用する。羽根車について、系が受け取る仕事は
おもりの位置エネルギーの減少 $mgh$ に等しい。導線を通って境界を横切る電気エネルギーは
電気的仕事として数える。
\end{indication}

\begin{solution}
\textbf{問1。} 温度差は $T_B-T_A=5{,}0$ K である。したがって
\[
\Delta U=C(T_B-T_A)
=2{,}00\times5{,}0
=10{,}0\ \mathrm{kJ}.
\]
この値は二つの平衡状態 $A$ と $B$ だけに依存する。

\textbf{問2。} 第一手順では仕事を受け取らないので $W_1=0$ である。
第一法則より
\[
Q_1=\Delta U=10{,}0\ \mathrm{kJ}.
\]
エネルギーが系に入るため、熱は正である。

第二手順では壁が断熱されているので $Q_2=0$ である。
内部エネルギー変化はすべて羽根車の力学的仕事による：
\[
W_2=\Delta U=10{,}0\ \mathrm{kJ}.
\]
$W_2=mgh$ だから
\[
m=\frac{W_2}{gh}
=\frac{10{,}0\times10^3}{9{,}81\times5{,}00}
\simeq 204\ \mathrm{kg}.
\]
この大きな桁は、わずかな温度上昇にもすでに相当量の力学的エネルギーが対応することを示す。

\textbf{問3。} 第一法則から
\[
W_3=\Delta U-Q_3=10{,}0-4{,}00=6{,}00\ \mathrm{kJ}.
\]
この仕事は正である。電源が系へエネルギーを与えるからである。
$W_3=\mathcal P\,\Delta t$ より
\[
\Delta t=\frac{W_3}{\mathcal P}
=\frac{6{,}00\times10^3}{300}
=20{,}0\ \mathrm{s}.
\]

\textbf{問4。} 三つの経路はそれぞれ
\[
(Q_1,W_1)=(10{,}0,0)\ \mathrm{kJ},
\qquad
(Q_2,W_2)=(0,10{,}0)\ \mathrm{kJ},
\]
および
\[
(Q_3,W_3)=(4{,}00,6{,}00)\ \mathrm{kJ}
\]
を与える。いずれも $Q+W$ は $10{,}0$ kJ で、同じ変化 $\Delta U$ を生む。
内部エネルギーは状態関数だが、熱と仕事は過程中のエネルギー移動を表す。
最終状態 $B$ で系は内部エネルギーをもつが、「熱」も「仕事」も含まない。
\end{solution}
\end{exo}

\begin{exo}[一定外圧下の圧縮]{compression-pression-exterieure-constante}
\seoready{true}
\keywords{圧力による仕事, 一定外圧, 圧縮, 膨張, 真空膨張}

気体を一定外圧 $P_0$ のもとで体積 $V_A$ から体積 $V_B<V_A$ まで圧縮する。

\begin{enumerate}
\item 気体が受け取る仕事を求めよ。
\item 気体が同じ外圧 $P_0$ に逆らって $V_B$ から $V_A$ に戻る。受け取る仕事を求め、圧縮時の仕事と比較せよ。
\item 気体が真空中で $V_B$ から $V_A$ へ膨張する場合を改めて考えよ。得られた三つの符号を解釈せよ。
\end{enumerate}

\begin{indication}
$W=-\int P_{\mathrm{ext}}\,\mathrm{d}V$ を用いる。過程が準静的でない場合も含め、
積分すべきなのは\emph{外部}圧力である。
\end{indication}

\begin{solution}
\textbf{問1。} 外圧は一定なので
\[
W_{A\to B}=-P_0(V_B-V_A)=P_0(V_A-V_B)>0.
\]
圧縮時、気体は仕事を受け取る。

\textbf{問2。} 戻りの過程では
\[
W_{B\to A}=-P_0(V_A-V_B)<0.
\]
二つの過程は同じ一定外圧に逆らって行われるため、二つの仕事は互いに反対符号である。

\textbf{問3。} 真空中では $P_{\mathrm{ext}}=0$ なので
\[
W_{B\to A}^{\mathrm{vide}}=0.
\]
したがって、膨張が自動的に負の仕事を意味するわけではない。実際に外力を押し返す必要がある。
\end{solution}
\end{exo}
